In view of the considerable surplus with which the balance account closes, an increase in the pension for the coming period appears permissible, and I have the following observations to make about the size of the increase.

If the question is asked by how much the variable part of the widow's pension must be increased, with all other arrangements retained, to bring debit and credit in the balance account into exact equality, a simple calculation gives this increase \(=\) 13 Thl. 7 Ggr. If a smaller increase is chosen, the revised balance account still closes with a surplus; it closes with a deficit if a larger increase is adopted. Thus, if the variable part of the pension were set at 60 thalers from now on, so that each eligible widow would receive a total of 260 thalers as long as their number does not exceed 18, and otherwise, in addition to the fixed part of 200 thalers, her corresponding share of the total sum of 1080 thalers, the balance account, drawn up in the same way as above, would still close with a surplus of 1695 thalers. There is therefore not only full justification for this measure but also the prospect of a further increase after a few years, provided that no \textit{large} losses occur, the benefit of the higher interest rate continues, and, while the number of widows has not yet reached 18, something is left over each year for the time being from the annual sum allowed for in the calculation.

The balance account would, however, close with a deficit of 877 thalers if the increase were set at 15 thalers, and with a deficit of 3448 thalers if it were set at 20 thalers. A deficit as small as that in the former case would be offset within a short time for the reasons just given. Thus there would, in itself, be no cause for concern at all in setting the variable part of the pension at 65 thalers (and consequently, when there are more than 18 widows, distributing the sum of 1170 thalers). Perhaps, however, one would be reluctant to depart from the practice always observed until now, by which the half-yearly pension has always amounted to a whole number of pistoles (a small convenience that will, of course, disappear of itself as soon as the number of widows has risen above 18). But I would consider setting the variable part at 70 thalers at this stage inadvisable on \SourceLeaf{174}principle alone, although, in the present favourable circumstances, I readily share the hope that the deficit of 3448 thalers would be considerably reduced in the coming years.

\medskip
I believe that in the foregoing I have brought together sufficient material for the university church deputation to make a decision with informed confidence; but I am glad to provide further explanations if these should be considered necessary.

Göttingen, 19 October 1851.\hfill C. F. \textsc{Gauss}.
\begin{center}\rule{.28\linewidth}{.3pt}\end{center}
